\chapter{溶液里发生了什么}

\begin{kaichang}
在厨房里做个小动作：接一杯清水，舀一勺食盐倒进去，搅一搅。十秒钟后，盐不见了。杯子还是那杯透明的液体，你看不到任何颗粒——但尝一口，是咸的。

现在请你对三个问题先猜后看：
\begin{itemize}[itemsep=2pt]
\item 盐“化”进水里，这杯水变重了吗？还是盐把重量“弄丢”了？
\item 水面这一口和杯底那一口，咸度一样吗？放一个月呢？
\item 把这杯盐水放着，让水慢慢干掉，盐还会回来吗？
\end{itemize}

把猜测写在纸上。这一章读完，你会发现“盐不见了”大概是厨房里最误导人的一句话——在粒子的世界里，没有任何东西消失，只有一场规模惊人的重新排队。
\end{kaichang}

\section{先称一称：没有东西消失}

先别急着讲道理，做一件任何科学家都会先做的事：称量。

取 \ud{250}{g} 水放在电子秤上，记下读数；再倒进 \ud{5}{g} 盐，搅匀。读数是多少？\ud{255}{g}，一两不少。把瓶子密封，放一个月再称，还是 \ud{255}{g}。盐的重量一点没丢，只是秤“看不见”它在哪儿了。

再看几个能观察的事实：

\begin{itemize}[itemsep=2pt]
\item 搅好的盐水，尝水面和尝杯底，咸度一样。而且放多久都一样——盐不会像沙子那样沉底。把一把沙子搅进水，几分钟就分层；盐水放一年也不分层。
\item 把盐水倒进浅盘，让水慢慢蒸发，白色晶体重新出现。尝一尝，是盐；称一称，接近原来的 \ud{5}{g}。
\item 同样的戏法不只属于盐：糖溶进水里是甜的，感冒冲剂溶进温水里整杯水都变味，汽水里溶着看不见却喝得出来的二氧化碳。碘酒则是碘溶在酒精里——\textbf{溶剂不一定是水}，只是水最常见。
\end{itemize}

“溶质分布均匀、放多久也不分层不沉淀”——满足这两条，就是一杯\textbf{溶液}。被溶的叫\textbf{溶质}（盐、糖、二氧化碳），负责溶的叫\textbf{溶剂}（水、酒精）。

所以“溶解”这个词的真实含义是：\textbf{溶质的粒子被拆开、打散，均匀地掺进溶剂的粒子里去了}。它还在，只是小到你看不见。小到什么程度？该把镜头推进去了。

\section{镜头推进：水分子拆围墙}

回忆两笔旧账。第 18 章说过，食盐里根本没有一个个独立的“\chem{NaCl} 分子”，它是钠离子 \chem{Na^+} 和氯离子 \chem{Cl^-} 像棋盘格一样交替堆成的巨大晶体，靠正负电荷的吸引力粘在一起。第 23 章说过，水分子是弯的，氧原子那一头偏负电，两个氢原子那一头偏正电——它是个小小的“电磁铁”。

现在把盐晶体放进水里，水面上看不见的攻防战开始了。晶格表面的 \chem{Na^+} 一露头，周围水分子的负端（氧端）立刻围上去；表面的 \chem{Cl^-} 则被水分子的正端（氢端）围住。几十个水分子七手八脚地拽一个离子，拽下来之后还层层围住不放，像给离子套上了一个“水分子外壳”。这个包围过程叫\textbf{溶剂化}；溶剂是水时，特称\textbf{水合}。被水合外壳裹住的离子，在水分子永不停歇的推挤下随机游走，从杯底一路扩散到水面——这就是你尝到“每口一样咸”的原因。

糖走的是另一条路。蔗糖是电中性的分子，拆不成离子；水分子干脆一拥而上，把整个蔗糖分子包围、带走。分子本身完好无损地漂在水里，所以糖水还是甜的——甜是蔗糖分子的性质，分子没坏，甜味就没坏。

注意一个关键词：\textbf{一直在动}。溶解不是把盐“放进”水里就结束了，而是离子和分子从此混入人群，跟着水分子日夜不停地乱撞。一杯看似安静的盐水，内部是一场永不散场的拥挤舞会。

规模有多大？算一笔账。食盐的摩尔质量约 \ud{58.5}{g/mol}，\ud{1}{g} 盐大约是 $1/58.5\approx 0.017$ mol。每摩尔含 $\ud{6.02\times 10^{23}}{}$ 个“\chem{NaCl} 单元”，每个单元入水后拆成 2 个离子，所以 \ud{1}{g} 盐入水变成约 $0.017\times 6.02\times 10^{23}\times 2\approx 2\times 10^{22}$ 个离子。一小勺盐撒进一杯水，平均每立方毫米分到大约 $10^{17}$ 个离子——比全地球海滩上的沙子总数还多得多。你的眼睛对这么小的粒子毫无办法，难怪你以为它“消失了”。

\section{溶解有上限：饱和、过饱和和气泡}

把水分子想象成搬运工，那它们的工作量有没有上限？有，而且你早就见过：往水里不断加盐，搅，再加，再搅，总有一刻，新加的盐沉在杯底，怎么搅都不再变少。这杯溶液\textbf{饱和}了。

但“饱和”绝不是“停止”。晶格表面的离子仍在被水分子撬走，同时水里乱撞的水合离子也会撞上晶体表面、重新卡进格子——\textbf{溶解和结晶同时在发生}。刚开时“出走”的多、“回家”的少，盐净减少；等水里离子够多，“回家”的速度追上了“出走”的速度，进出恰好相抵。表面看起来风平浪静，粒子层面却是两股方向相反的洪流在对冲。这种“动着的平衡”在化学里无处不在，第 28 章我们会专门对付它。

这里藏着一个经典证据：把一粒形状不规则的盐粒用细线吊进它自己的饱和盐水里，几天后取出——盐粒的重量几乎没变（进出相抵），但棱角变圆滑了（表面的粒子被换了一遍）。如果饱和意味着“什么都不发生”，这粒盐就该原封不动。

饱和线不是固定的，它听两个旋钮的指挥：

\begin{itemize}[itemsep=2pt]
\item \textbf{温度}。对大多数固体溶质，水越热，溶得越多。以硝酸钾 \chem{KNO_3} 为例：$20\,^{\circ}\mathrm{C}$ 时 \ud{100}{g} 水大约溶 \ud{32}{g}，加热到 $80\,^{\circ}\mathrm{C}$ 能溶 \ud{169}{g}——五倍多。反过来，把热乎乎溶得满满的溶液冷却下来，多出来的溶质就得“让座”，重新析出晶体。蜂蜜放久了结出糖粒、冰糖在糖浆里慢慢长大，都是这个原理。
\item \textbf{压强}。对固体和液体溶质，压强几乎无所谓；但对\textbf{气体}溶质，压强说了算。汽水是在高压下把二氧化碳 \chem{CO_2} 硬压进水里的，你一拧开瓶盖，压强骤降，溶不住的气体立刻争先恐后地逃逸——那就是你看到的满杯气泡。
\end{itemize}

温度对气体的规矩恰好相反：水越热，溶的气体越\textbf{少}。夏天闷热的午后，鱼塘里的鱼浮到水面“喘气”，就是因为热水里氧气不够。一杯冷水在室温下放久了内壁冒出小气泡，也是原本溶解的空气在“退房”。

还有一种狡猾的状态：把热的浓溶液小心地、毫无振动地冷却，溶质有时来不及析出，整杯溶液处于一种“超编却暂时没人查岗”的状态——\textbf{过饱和}。它极不稳定：丢进一粒小晶体，甚至只是晃一下，过量溶质就会在几秒内沿着晶种疯狂结晶，整杯液体瞬间“长”出晶体森林。冬天用的可反复掰铁片发热的暖手宝，里面正是一袋过饱和醋酸钠溶液；至于它结晶时为什么会发热，那是第 27 章能量账本的题目。

\section{糖水点不亮灯泡，盐水可以}

下面这个对比，是理解溶液导电的全部关键。

想象一个简单的电路：电池、小灯泡，两根导线末端露出金属电极。把电极插进不同的液体：纯净水，灯泡几乎不亮；浓糖水，还是不亮；换成盐水，灯泡亮了；换成白醋，灯丝泛起微弱的光。

电流的本质是带电粒子的定向移动。固体金属里，跑腿的是电子；液体里没有自由电子，想导电就得靠\textbf{能自由移动的带电粒子}。纯净水分子整体电中性，糖分子也电中性——糖水是整整一杯中性分子，谁也驮不动电荷。而盐水里游荡着大量 \chem{Na^+} 和 \chem{Cl^-}，一通电，正离子奔向负极、负离子奔向正极，两股离子流就把电流“接力”传了过去。

溶于水（或熔融）后能导电的物质叫\textbf{电解质}，不导电的叫\textbf{非电解质}。盐、酸、碱大多是电解质；糖、酒精是非电解质。注意一个配套事实：固态食盐不导电——离子被锁死在晶格里，动弹不得（第 18 章）。同一把盐，固体不导电、溶于水导电，差别只在“离子能不能自由移动”。

说到这儿，符号终于可以登场了。我们用 (s) 表示固态、(aq) 表示“水合状态”（拉丁语 aqua，水），把盐入水写成：

\[\chem{NaCl(s)} \rightarrow \chem{Na^+(aq)} + \chem{Cl^-(aq)}\]

一个 \chem{NaCl} 单元变成一个水合钠离子和一个水合氯离子。而蔗糖的“方程式”就朴素得多：

\[\chem{C_{12}H_{22}O_{11}(s)} \rightarrow \chem{C_{12}H_{22}O_{11}(aq)}\]

分子原封不动，只是从固体搬进了水里。两个式子摆在一起，“盐水电、糖水不电”的原因一目了然：左边都还是分子或晶格，右边一个拆出了带电粒子、一个没有。

这件事不只是课堂把戏。你的汗水是货真价实的电解质溶液，所以湿手摸开关远比干手危险——水膜里的离子帮电流开路。医院的生理盐水、养鱼的换水、锅炉的水垢处理，全都绕着“水里有哪些离子、有多少”打转。

\section{科学家怎样知道：水里的盐真的拆成了两半}

\begin{zhengju}
今天你觉得“盐在水里变成离子”天经地义，但在 1880 年代，这是个近乎荒唐的主张。当时的主流看法是：化合物只有在通电时才会被电流“劈”成离子——电是因，离子是果。

麻烦先从荷兰化学家范特霍夫那里冒出来。他在 1885 年前后系统测量了溶液的\textbf{渗透压}（后面细讲，你可以先把它理解成“溶液渴水的程度”），发现一个漂亮的规律：糖溶液的渗透压和“溶质粒子越多、压强越大”的公式完美吻合，仿佛溶质粒子像气体分子一样各自独立行动。可是换成食盐，同样浓度的溶液，渗透压几乎是公式的\textbf{两倍}；换成氯化钙 \chem{CaCl_2}，接近\textbf{三倍}。为什么盐和糖不守同一个规矩？

瑞典的年轻博士生阿伦尼乌斯在 1884 年的博士论文里给出了一个大胆解释：电解质一入水就\textbf{自动}解离成带电离子，不需要任何电流。如果一份 \chem{NaCl} 变成 2 个独立粒子，渗透压自然是 2 倍；\chem{CaCl_2} 变成 3 个，自然是 3 倍——范特霍夫的怪数据立刻说得通了。这个比值后来被叫作“范特霍夫因子”。

当时的教授们不买账，答辩只给了勉强及格的低分。反对理由听起来相当有力：金属钠遇水会爆炸，氯气是剧毒气体，你说水里安安静静地漂着钠离子和氯离子？阿伦尼乌斯的回答朴素而深刻：\textbf{离子和原子是两种不同的东西}。钠原子有一颗急着脱手的电子所以暴躁，钠离子已经交出了电子所以温顺（第 10 章的电子楼层逻辑）；水里的 \chem{Na^+} 跟金属钠的关系，就像退休老人和拳击手——同一个人，完全不同的状态。

争论持续了很多年，直到越来越多的独立证据汇聚过来：导电性、冰点变化、X 射线晶体学直接拍出离子晶体，全都指向同一结论。1903 年，阿伦尼乌斯拿到了诺贝尔化学奖。故事还有后半截：人们很快发现，浓溶液里范特霍夫因子会缩水——\chem{NaCl} 的“2 倍”在浓溶液里跌到 1.9、1.8。这不是电离学说错了，而是离子多了会互相牵绊、结伴行动。1923 年德拜和休克尔把这一层修正算了进去。科学理论很少一锤定音，它是这样一层层补强的。
\end{zhengju}

\section{既不算溶，也不算不溶：胶体和牛奶}

现在来为难一下我们的分类。牛奶看起来均匀、放几天也不沉淀，它是溶液吗？

用一束光就能审问它。黑暗里拿激光笔（或手机手电筒）分别照射一杯盐水和一杯兑了水的牛奶：盐水里几乎看不到光的路，牛奶里却出现一条明亮的“光柱”。这是英国物理学家丁达尔发现的现象：光柱的出现，说明液体里有\textbf{大到足以散射光的颗粒}。盐水的离子直径不到 1 纳米，光波绕过去毫无感觉；牛奶里的脂肪球和蛋白质团块直径在几百纳米，正好把光撞得四散。

按颗粒大小，分散系排成一条连续的谱系：

\begin{itemize}[itemsep=2pt]
\item \textbf{真溶液}（颗粒小于约 1 纳米）：离子、小分子级别。透明，永不沉降。盐水、糖水。
\item \textbf{胶体}（约 1 纳米到 1 微米）：颗粒小到不会沉降，却大到能散射光。牛奶、豆浆、血液、果冻、雾和云。
\item \textbf{悬浊液、乳浊液}（大于约 1 微米）：颗粒迟早沉降或分层。泥水、没摇匀的油水混合物。
\end{itemize}

牛奶还有个身份：它是\textbf{乳浊液}被“稳住”了的版本。油和水本来势不两立（第 23 章讲过，极性分子和非极性分子互相回避），混一混也会很快分层。但牛奶里的蛋白质分子一头亲水、一头亲油，它们包住油滴，亲水的头朝外面对水，油滴就被“伪装”成水喜欢的样子，稳定地悬浮着。这种物质叫\textbf{乳化剂}。蛋黄酱靠蛋黄里的卵磷脂乳化，洗洁精去油污靠的也是这一手——它的分子同样一头亲水一头亲油，把油滴从碗上“撬”下来包进水里冲走。

胶体不是教科书里的冷门名词：你血管里流淌的血浆、豆腐的凝固、蓝天的颜色、自来水的絮凝净化，全是胶体化学。

\section{一层只放行一部分粒子的膜：渗透}

最后一站，是全章最重要的概念。它今天解释一棵腌黄瓜，明天解释你身体里的每一个细胞。

有些薄膜有一种选择性：让小个子的水分子穿过，却把大个子的溶质粒子挡在门外。这种膜叫\textbf{半透膜}。鸡蛋壳内层那层白色的蛋膜、香肠的肠衣、实验室的玻璃纸，都是半透膜。

做一个思想实验：用半透膜把一个槽子隔成两边，一边纯水，一边浓糖水。几小时后你会看到：糖水那边的液面升高了，纯水那边降低了。\textbf{水穿过膜，净流向了糖水一侧}。这个现象叫\textbf{渗透}。

水为什么“偏要去”浓的那边？先排除一个诱人的错误说法：浓溶液并没有伸出一只手把水“吸”过去。真相是纯统计的。膜两侧的水分子都在乱撞、都有机会钻过膜孔，但浓溶液那边，一部分撞膜的位置被溶质粒子占着，单位时间里从浓侧出发穿过膜的水分子就比稀侧少。两边都在过，稀侧过得多、浓侧过得少，净效果就是水涌向浓侧——直到升高的液柱产生足够的压强，把多出来的水“压回去”，进出重新相抵。平衡时这股压强差，就是前面提过的\textbf{渗透压}。

\begin{qiaoliang}
渗透压有多大？范特霍夫发现它和气体压强服从几乎同一个公式：$\pi = iMRT$，其中 $M$ 是溶质的物质的量浓度，$i$ 是“一份溶质变成几个粒子”的范特霍夫因子，$R$ 是常数，$T$ 是绝对温度。拿医院输液用的生理盐水（$0.9\%$ 的食盐水）算一算：每升含盐 \ud{9}{g}，即 $9/58.5\approx 0.154$ mol；食盐入水拆成 2 个离子，$i\approx 2$；代入 $T\approx 310$ K（体温），$R\approx \ud{0.082}{L\cdot atm/(mol\cdot K)}$，得

\[\pi \approx 2\times 0.154\times 0.082\times 310 \approx 7.8\ \mathrm{atm}\]

约 7.8 个大气压——相当于一根 80 米高的水柱的压力！这么小半勺盐溶进一瓶水，居然能“调动”如此巨大的力量，难怪渗透压绝不能怠慢。而 7.8 个大气压恰好与你体液的渗透压相等，这就是输液要用 $0.9\%$ 的原因：\textbf{等渗}，不多不少。换成纯水直接进血管，红细胞会吸水胀破；换成浓盐水，细胞又会失水皱缩。剂量差之毫厘，细胞命悬一线。
\end{qiaoliang}

渗透早已渗透进你的生活：盐撒在黄瓜上，黄瓜细胞里的水渗出来，成了腌菜汁；糖拌西红柿的盘底那层甜水，是同一个戏法；海鱼喝海水、靠鳃排盐，淡水鱼则几乎不喝水、拼命排尿——两种鱼在淡水和海水里各自进化出对付渗透的策略，互换环境都是死路。

人类还把这个原理倒过来用：对海水施加\textbf{超过}渗透压的压力，水就被迫反着穿过半透膜，把盐留在身后——这就是\textbf{反渗透}海水淡化。以色列、新加坡和波斯湾沿岸城市的大量饮用水，都是这么从海里“挤”出来的。你的净水器里多半也有一支反渗透膜滤芯。

\begin{shiyan}
\textbf{会喝水、会失水的“裸蛋”}

材料：生鸡蛋 2 枚、白醋 1 瓶、浓糖水（或蜂蜜）一杯、清水一杯、透明玻璃杯 3 个、保鲜膜、厨房秤（可选）。

步骤：
\begin{enumerate}[itemsep=2pt]
\item 把鸡蛋轻轻放入杯中，倒入白醋没过鸡蛋，盖上保鲜膜。蛋壳表面会立刻冒小气泡——那是醋酸在溶解蛋壳里的碳酸钙，放出的正是 \chem{CO_2}。
\item 静置 24～48 小时，中途换一次醋。等蛋壳完全消失，就得到只剩一层蛋膜包裹的“裸蛋”，半透明、有弹性。这层蛋膜就是天然半透膜：水能过，糖和蛋白质不能过。
\item 裸蛋 A 泡进浓糖水，裸蛋 B 泡进清水，各静置 12 小时，中途可翻个面。
\end{enumerate}

观察点：A 蛋明显变皱、变小、变轻（如果有秤，称一称前后质量）；B 蛋胀大、变沉、摸起来紧绷。为什么？蛋内液体的浓度比清水高、比浓糖水低——水总是净流向更浓的一侧。

安全提示：白醋有刺激性气味，操作时保持通风，溅到眼睛立即用大量清水冲洗；实验后的鸡蛋已被醋和外界微生物污染，\textbf{不能吃}；玻璃杯轻拿轻放。
\end{shiyan}

\section{这套图景哪里会失灵}

本章的模型干净好用，但先说好它的边界，免得你将来被它骗：

\begin{itemize}[itemsep=2pt]
\item “浓度越高渗透压越大”“范特霍夫因子等于离子个数”都只在\textbf{稀溶液}里精确成立。浓溶液里离子互相干扰、水分子不够分配，一切都开始偏离简单公式——阿伦尼乌斯的理论当年就靠这层修正才站稳。
\item “溶解度随温度升高而增大”对多数固体成立，却有反例：氢氧化钙（熟石灰）反而越热越难溶，气体更是集体反着来。经验规律不是铁律，用之前先查数据。
\item “溶解”与“化学反应”之间没有清晰的围墙。食盐入水，离子本身没变，我们说这是物理过程；但 \chem{CO_2} 入水有一部分会变成碳酸，金属钠入水干脆炸给你看。从“纯掺和”到“彻底变身”是一条连续光谱，下一篇（第 26 章起）我们正式处理“变身”的那一半。
\item 胶体和真溶液的界限（1 纳米）是人为划的。纳米颗粒、超大蛋白质、病毒，都骑在这条线上——分类是工具，不是自然界的裂缝。
\end{itemize}

\begin{wuqu}
“渗透是水被高浓度‘吸引’过去的。”——这是拟人化惹的祸。分子没有眼睛，看不到对面有多浓；也没有任何跨越膜的“吸力”。每个水分子只管漫无目的地乱撞，撞到哪边算哪边。净方向来自一个冷酷的统计事实：浓的那边，撞膜的水分子里有一部分被溶质粒子“顶替”了。验证方法很简单：把 U 形管横过来、倒过来，渗透照样发生，方向和重力毫无关系——“吸力说”没法解释这一点，统计说可以。记住这本书的口头禅：别问粒子“想”去哪，只问哪条路走的人多。
\end{wuqu}

\section{那勺盐到底去哪了}

回到开场的三个猜想，对答案。

第一，盐水当然变重了：水的重量加盐的重量，一分不少。质量守恒的底气来自原子守恒——溶解从头到尾没有创造或消灭任何一个原子，只是把它们从整齐的晶格队列改编成了游荡的水合离子大军（这背后的普遍法则，第 26 章会立成铁律）。

第二，每口一样咸，而且永远一样咸。约 $2\times 10^{22}$ 个离子在热运动的驱动下均匀散步，没有沉降、没有分层，也不打算回到杯底。

第三，把水蒸干，离子们失去水分子的包围，只能重新按正负交替的棋盘格排队——盐回来了，还是那把盐。

“消失”从来不是粒子的命运，只是你眼睛的极限。

\section{从一杯盐水，到一座被膜围住的城市}

这一章其实一直在为一个更大的故事铺路，你可能还没察觉。

你的血液是溶液，汗是溶液，眼泪是溶液，细胞里那汪细胞质也是溶液——一个成年人的身体，就是一套约 40 升、成分精确调控的水溶液系统。葡萄糖溶在血里奔赴全身，钠离子和钾离子在神经里掀起电信号，肾脏日夜不停地调节着你这“一杯溶液”的浓度。

而这一切的前提是：每一个细胞都被一层膜包围着。那层膜首先就是一种半透膜——水能过，许多溶质不能随心所欲地过。于是我们本章的所有问题立刻变成了细胞的生死问题：渗透压失控，细胞会胀破或干瘪；离子浓度失衡，神经和心脏就停摆。但细胞膜远比鸡蛋膜高明：它在膜上装了只能单向通行的“水泵”、识别特定离子的“闸门”，把被动的渗透变成了主动的管理。一层没有生命的膜，怎样做出近乎“会思考”的选择？那是第二册第 51 章《细胞膜》的故事——在那里，你会看到本章的每一个概念（水合离子、浓度、半透膜、渗透压）原封不动地重新登场，只是舞台从烧杯换成了你自己。

\begin{tiaozhan}
溶质粒子还有一个统称“依数性”的本事：只数个数、不问身份地改变溶剂的性质——除了渗透压，还包括\textbf{凝固点下降}和\textbf{沸点升高}。稀溶液里，凝固点下降 $\Delta T = i\cdot K_f\cdot m$，其中 $K_f$ 是溶剂的常数（水为 $1.86\,^{\circ}\mathrm{C\cdot kg/mol}$），$m$ 是质量摩尔浓度。算一笔冬天的账：把 1 mol 食盐（\ud{58.5}{g}）撒进 \ud{1}{kg} 冰雪融成的水里，$i\approx 2$，凝固点下降约 $2\times 1.86\approx 3.7\,^{\circ}\mathrm{C}$——这就是市政往路面撒盐融雪的原理。但注意边界：气温低于约 $-10\,^{\circ}\mathrm{C}$ 时撒再多盐也无济于事，所以严寒地区得换氯化钙（$i\approx 3$，还放热）或干脆上铲雪车。反过来用：冰棍模具里的糖水比纯水更难冻，做实验测一测你家的冰箱冷冻室能不能冻硬饱和盐水？
\end{tiaozhan}

\section*{练一练}

\begin{sikaoti}
\item 画一幅“一滴盐水放大一亿倍”的粒子图：画出几个 \chem{Na^+}、\chem{Cl^-} 和它们周围的水分子，注意标出每个水分子氧端（负端）和氢端（正端）分别朝向哪种离子。图旁注明：这是人为的表示法，水合外壳时刻在解体和重组。
\item 把一粒表面粗糙的冰糖用线吊进它自己的饱和糖水里，密封放一周。预测它的质量和形状各自会怎样变化，并用“溶解与结晶双向对冲”解释。这个实验为什么能证明饱和不是“静止”？
\item 画一张物质流图：半透膜左侧是纯水，右侧是浓盐水，用箭头画出两侧水分子的穿越情况（箭头粗细表示流量大小），标出净流向和最终平衡状态。再把装置换成“左侧稀盐水、右侧浓盐水”，净流向变不变？
\item 医生给脱水病人输液时，为什么用 $0.9\%$ 的盐水而不是纯水或 $9\%$ 的盐水？请从红细胞（可以理解为水球外包一层半透膜）的角度，分别画出两种错误选择的后果示意图。
\item 给你一杯清水、一杯稀牛奶、一杯盐水和一支激光笔，设计一个实验区分三杯液体，并预测每杯的观察结果。
\item 估算：一瓶 \ud{500}{mL} 生理盐水（$0.9\%$）里大约有多少摩尔的离子？这些离子大约有多少个？（提示：先算溶质质量，再用范特霍夫因子。）
\end{sikaoti}
